%% LyX 2.5.1 created this file.  For more info, see https://www.lyx.org/.
%% Do not edit unless you really know what you are doing.
\documentclass[twoside,onecolumn,american,compsoc]{IEEEtran}
\usepackage[T1]{fontenc}
\usepackage[utf8]{inputenc}
\usepackage{mathtools}
\usepackage{amsmath}
\usepackage{amssymb}

\makeatletter
%%%%%%%%%%%%%%%%%%%%%%%%%%%%%% Textclass specific LaTeX commands.
\numberwithin{equation}{section}
\numberwithin{figure}{section}

%%%%%%%%%%%%%%%%%%%%%%%%%%%%%% User specified LaTeX commands.
\usepackage[left=1.4cm, right=1.5cm, top=1.5cm, bottom=1.5cm]{geometry}

\makeatother

\usepackage{babel}
\begin{document}

\section{Unconditional Linear Denoiser with Circulant Matrix}

A linear denoiser is $\boldsymbol{D}(\boldsymbol{y};\sigma)=\boldsymbol{W}_{\sigma}\boldsymbol{y}+\boldsymbol{b}_{\sigma}$,
but now consider the case that $\boldsymbol{W}_{\sigma}$ is constrained
to be circulant.
\begin{align*}
\mathcal{L}^{\text{circ}}_{\sigma} & =\mathbb{E}_{\boldsymbol{x_{0}}\sim p_{0},\boldsymbol{Z}\sim\mathcal{N}(0,\boldsymbol{I})}\|\boldsymbol{W}^{\text{circ}}_{\sigma}(\boldsymbol{x_{0}}+\sigma\boldsymbol{Z})+\boldsymbol{b}_{\sigma}-\boldsymbol{x_{0}}\|^{2}_{2}
\end{align*}

where $\boldsymbol{W}^{\text{circ}}_{\sigma}\in\mathbb{R}^{d\times d},\boldsymbol{x_{0}}\in\mathbb{R}^{d\times1}$.
Diagonalize $\boldsymbol{W}^{\text{circ}}_{\sigma}\coloneqq\boldsymbol{F}^{*}\text{diag}(\hat{\boldsymbol{W}})\boldsymbol{F}$,
where $\boldsymbol{F}$ is the unitary (across $d$) Discrete Fourier
Transform matrix $\boldsymbol{F}$ where $\boldsymbol{F}_{jk}=\frac{1}{\sqrt{d}}e^{-2\pi ijk/d}$
and $\boldsymbol{F}^{*}$ is the conjugate transpose where $\boldsymbol{F}^{*}_{jk}=\frac{1}{\sqrt{d}}e^{2\pi ijk/d}$,
and $\text{diag}(\hat{\boldsymbol{W}})$ is the Discrete Fourier Transform
of the first row of the circulant matrix $\boldsymbol{W}$, i.e. $\text{diag}(\hat{\boldsymbol{W}})_{j}\coloneqq\hat{w}_{j}=\sum^{d}_{k=0}w_{k}e^{-2\pi ijk/d}$
($w_{k}$ is the $k$-th entry of the first row of $\boldsymbol{W}$).
For $\boldsymbol{W}$ to have real entries, $\hat{w}_{j}=\overline{\hat{w}_{d-j}}$.

Recall that the optimal $\boldsymbol{b}_{\sigma}$ is $\boldsymbol{b}^{*}_{\sigma}=(\boldsymbol{\boldsymbol{I}-W}_{\sigma})\mathbb{E}\big[\boldsymbol{x_{0}}\big]$.
So we have:
\[
\mathcal{L}^{\text{circ}}_{\sigma}=\mathbb{E}_{\boldsymbol{x_{0}}\sim p_{0},\boldsymbol{Z}\sim\mathcal{N}(0,\boldsymbol{I})}\|\boldsymbol{W}^{\text{circ}}_{\sigma}(\boldsymbol{x_{0}}-\mathbb{E}\big[\boldsymbol{x_{0}}\big])+\sigma\boldsymbol{Z})-(\boldsymbol{x_{0}}-\mathbb{E}\big[\boldsymbol{x_{0}}\big])\|^{2}_{2}
\]

Let $\tilde{\boldsymbol{x_{0}}}=\boldsymbol{x_{0}}-\mathbb{E}\big[\boldsymbol{x_{0}}\big]$.
Then:
\begin{align*}
\mathcal{L}^{\text{circ}}_{\sigma} & =\mathbb{E}_{\boldsymbol{x_{0}}\sim p_{0},\boldsymbol{Z}\sim\mathcal{N}(0,\boldsymbol{I})}\|\boldsymbol{W}^{\text{circ}}_{\sigma}(\tilde{\boldsymbol{x_{0}}}+\sigma\boldsymbol{Z})-\tilde{\boldsymbol{x_{0}}}\|^{2}_{2}\\
 & =\text{Tr}\big(\boldsymbol{W}^{\text{circ}}_{\sigma}(\boldsymbol{\Sigma}_{\boldsymbol{p_{0}}}+\sigma^{2}\boldsymbol{I})\boldsymbol{W}^{\text{circ}\top}_{\sigma}-2\boldsymbol{W}^{\text{circ}}_{\sigma}\boldsymbol{\Sigma}_{\boldsymbol{p_{0}}}+\boldsymbol{\Sigma}_{\boldsymbol{p_{0}}}\big)
\end{align*}

Let $\boldsymbol{P}=\boldsymbol{F}\boldsymbol{\Sigma}_{\boldsymbol{p_{0}}}\boldsymbol{F}^{*}$,
then $\boldsymbol{P}$ is Hermitian, and $\boldsymbol{P}_{kk}=f^{H}_{k}\boldsymbol{\Sigma}_{\boldsymbol{p_{0}}}f_{k}$
where $f_{k}$ is the $k$-th column of $\boldsymbol{F}^{*}$, or
the conjugate of the $k$-th row of $\boldsymbol{F}$. Note that $\boldsymbol{P}_{kk}=\boldsymbol{P}_{d-k,d-k}$.
Since $\boldsymbol{W}^{\text{circ}}_{\sigma}$ is real valued, $\boldsymbol{W}^{\text{circ}\top}_{\sigma}=\boldsymbol{W}^{\text{circ}H}_{\sigma}$
($\boldsymbol{A}{}^{H}$ is the conjugate transpose matrix of $\boldsymbol{A}$).
So $\boldsymbol{W}^{\text{circ}\top}_{\sigma}=(\boldsymbol{F}^{*}\text{diag}(\hat{\boldsymbol{W}})\boldsymbol{F})^{H}=\boldsymbol{F}^{H}\text{diag}(\hat{\boldsymbol{W}})^{H}(\boldsymbol{F}^{*})^{H}=\boldsymbol{F}^{*}\text{diag}(\hat{\boldsymbol{W}})^{H}\boldsymbol{F}$.

Since $\boldsymbol{F}$ and $\boldsymbol{F}^{*}$ are both symmetric,
$\boldsymbol{W}^{\text{circ}\top}_{\sigma}=\boldsymbol{F}^{*}\text{diag}(\hat{\boldsymbol{W}})^{H}\boldsymbol{F}$
is also symmetric . Consider the first term:
\begin{align*}
 & \text{Tr}(\boldsymbol{F}^{*}\text{diag}(\hat{\boldsymbol{W}})\boldsymbol{F}(\boldsymbol{\Sigma}_{\boldsymbol{p_{0}}}+\sigma^{2}\boldsymbol{I})\boldsymbol{F}^{*}\text{diag}(\hat{\boldsymbol{W}})^{H}\boldsymbol{F})\\
 & =\text{Tr}(\boldsymbol{F}\boldsymbol{F}^{*}\text{diag}(\hat{\boldsymbol{W}})\boldsymbol{F}(\boldsymbol{\Sigma}_{\boldsymbol{p_{0}}}+\sigma^{2}\boldsymbol{I})\boldsymbol{F}^{*}\text{diag}(\hat{\boldsymbol{W}})^{H})\\
 & =\text{Tr}(\text{diag}(\hat{\boldsymbol{W}})(\boldsymbol{P}\text{+\ensuremath{\sigma^{2}\boldsymbol{I}})diag}(\hat{\boldsymbol{W}})^{H})\\
 & =\sum^{d}_{k=1}\|\hat{w}_{k}\|^{2}(\boldsymbol{P}_{kk}+\sigma^{2})
\end{align*}
Consider the second term:
\begin{align*}
 & \text{Tr}(\boldsymbol{F}^{*}\text{diag}(\hat{\boldsymbol{W}})\boldsymbol{F}\boldsymbol{\Sigma}_{\boldsymbol{p_{0}}})\\
 & =\text{Tr}(\text{diag}(\hat{\boldsymbol{W}})\boldsymbol{F}\boldsymbol{\Sigma}_{\boldsymbol{p_{0}}}\boldsymbol{F}^{*})\\
 & =\text{Tr}(\text{diag}(\hat{\boldsymbol{W}})\boldsymbol{P})=\sum^{d}_{k=1}\hat{w}_{k}\boldsymbol{P}_{kk}
\end{align*}
Since $\boldsymbol{P}_{kk}=\boldsymbol{P}_{d-k,d-k}$, pairing $k$
and $d-k$ in this sum gives us: $\hat{w}_{k}\boldsymbol{P}_{kk}+\hat{w}_{d-k}\boldsymbol{P}_{d-k,d-k}=\boldsymbol{P}_{kk}(\hat{w}_{k}+\bar{\hat{w}}_{k})=\boldsymbol{P}_{kk}2\text{Re}(\hat{w}_{k})$.
So $\text{Tr}(\boldsymbol{F}^{*}\text{diag}(\hat{\boldsymbol{W}})\boldsymbol{F}\boldsymbol{\Sigma}_{\boldsymbol{p_{0}}})=\sum^{d}_{k=1}\text{Re}(\hat{w}_{k})\boldsymbol{P}_{kk}$.
Then:
\[
\mathcal{L}^{\text{circ}}_{\sigma}=\sum^{d}_{k=1}(\|\hat{w}_{k}\|^{2}(\boldsymbol{P}_{kk}+\sigma^{2})-2\text{Re}(\hat{w}_{k})\boldsymbol{P}_{kk})+\text{Tr}(\boldsymbol{\Sigma}_{\boldsymbol{p_{0}}})
\]
Now observe that each $\hat{w}_{k}$ is independent and so we can
optimize for each $\hat{w}_{k}$ independently. Let each $\hat{w}_{k}=a+bi$
where $a,b\in\mathbb{R}$. Our objective is to minimize:
\[
\|\hat{w}_{k}\|^{2}(\boldsymbol{P}_{kk}+\sigma^{2})-2\text{Re}(\hat{w}_{k})\boldsymbol{P}_{kk}=(a^{2}+b^{2})(\boldsymbol{P}_{kk}+\sigma^{2})-2a\boldsymbol{P}_{kk}
\]
Note that $\boldsymbol{P}_{kk}$ is real and non-negative, so to achieve
the minimum of $\mathcal{L}^{\text{circ}}_{\sigma}$, it ought to
be that $b=0$ and solving for the minimum of w.r.t $a$:
\[
\nabla_{a}a^{2}(\boldsymbol{P}_{kk}+\sigma^{2})-2a\boldsymbol{P}_{kk}=2a(\boldsymbol{P}_{kk}+\sigma^{2})-2\boldsymbol{P}_{kk}=0
\]
so:
\[
a=\frac{\boldsymbol{P}_{kk}}{\boldsymbol{P}_{kk}+\sigma^{2}}=\hat{w}^{*}_{k}
\]
Therefore:
\[
\mathcal{L}^{\text{circ}}_{\sigma}=\sum^{d}_{k=1}\frac{-\boldsymbol{P}^{2}_{kk}}{\boldsymbol{P}_{kk}+\sigma^{2}}+\text{Tr}(\boldsymbol{\Sigma}_{\boldsymbol{p_{0}}})
\]
Note that $\sum^{d}_{k=1}\boldsymbol{P}_{kk}=\sum^{d}_{k=1}\text{Tr}(f^{H}_{k}\boldsymbol{\Sigma}_{\boldsymbol{p_{0}}}f_{k})=\text{Tr}(\boldsymbol{\Sigma}_{\boldsymbol{p_{0}}}\sum^{d}_{k=1}f_{k}f^{H}_{k})=\text{Tr}(\boldsymbol{\Sigma}_{\boldsymbol{p_{0}}}\boldsymbol{F}\boldsymbol{F}^{*})=\text{Tr}(\boldsymbol{\Sigma}_{\boldsymbol{p_{0}}}\boldsymbol{I})=\text{Tr}(\boldsymbol{\Sigma}_{\boldsymbol{p_{0}}})$,
so:
\[
\mathcal{L}^{\text{circ}}_{\sigma}=\sum^{d}_{k=1}\frac{-\boldsymbol{P}^{2}_{kk}}{\boldsymbol{P}_{kk}+\sigma^{2}}-\boldsymbol{P}_{kk}=\sigma^{2}\sum^{d}_{k=1}\frac{\boldsymbol{P}_{kk}}{\boldsymbol{P}_{kk}+\sigma^{2}}
\]
Recall that the regular optimal denoiser loss is:
\[
\mathcal{L}_{\sigma}=\sigma^{2}\sum^{d}_{k=1}\frac{\lambda_{k}}{\lambda_{k}+\sigma^{2}}
\]
where $\lambda_{k}$ is the eigenvalues of $\boldsymbol{\Sigma}_{\boldsymbol{p_{0}}}$.
Since $\boldsymbol{F}\boldsymbol{F}^{*}=\boldsymbol{I}$, the transformation
$\boldsymbol{P}=\boldsymbol{F}\boldsymbol{\Sigma_{x_{0}}}\boldsymbol{F}^{*}$
preserves the eigenvalues of $\boldsymbol{\Sigma_{x_{0}}}$. By the
Schur-Horn Theorem, for any Hermitian matrix, the vector of its true
eigenvalues always majorizes the vector of its diagonal entries, i.e.
$\boldsymbol{\lambda}\succ\text{diag}(\boldsymbol{P})$. Let $f(x)=\frac{x}{x+\sigma^{2}}$;
it can be confirmed that $f(x)$ is concave since the second derivative
$f''(x)=\frac{-2\sigma^{2}}{(x+\sigma^{2})^{3}}$ is strictly negative
for $x\geq0$ (note that $\lambda_{k}$ and $\boldsymbol{P}_{kk}$
are non-negative). By the Karamata's Inequality/Majorization Theorem,
if a vector $\boldsymbol{x}$ majorizes $\boldsymbol{y}$ and $f$
is concave, then $\sum^{d}_{k=1}f(x_{k})\leq\sum^{d}_{k=1}f(y_{k})$.
So:
\[
\mathcal{L}^{\text{circ}}_{\sigma}=\sigma^{2}\sum^{d}_{k=1}\frac{\boldsymbol{P}_{kk}}{\boldsymbol{P}_{kk}+\sigma^{2}}\geq\sigma^{2}\sum^{d}_{k=1}\frac{\lambda_{k}}{\lambda_{k}+\sigma^{2}}=\mathcal{L}_{\sigma}
\]
Constraint for $\boldsymbol{W}_{\sigma}$ to be circulant doesn't
buy us anything in the linear denoiser case; since circulant matrices
is still a subset of matrices and so the optimal denoiser loss ought
to be the lower bound for the circulant matrix denoiser loss.

\section{Conditional Linear Denoiser with Circulant Matrix}

Now similarly consider the conditional denoiser case, where we also
constrain $\boldsymbol{W}_{\sigma}$ to be circulant. Recall the optimal
$\boldsymbol{b}_{\sigma}$ is $\boldsymbol{b}^{*}_{\sigma}=(\boldsymbol{\boldsymbol{I}-W}_{\sigma})\boldsymbol{\mu_{p_{0}}}-\boldsymbol{V}_{\sigma}\boldsymbol{\mu_{U}}$
and $\boldsymbol{V}^{*}_{\sigma}=(\boldsymbol{\boldsymbol{I}-W}_{\sigma})\text{Cov}(\boldsymbol{x_{0}},\boldsymbol{U})\boldsymbol{\Sigma}^{-1}_{\boldsymbol{U}}$
\begin{align*}
\mathcal{L}^{\text{circ}}_{\sigma,\boldsymbol{U}} & =\mathbb{E}_{\boldsymbol{x_{0}}\sim p_{0},\boldsymbol{Z}\sim\mathcal{N}(0,\boldsymbol{I}),\boldsymbol{U}}\|\boldsymbol{W}^{\text{circ}}_{\sigma}(\boldsymbol{x_{0}}+\sigma\boldsymbol{Z})+\boldsymbol{V}_{\sigma}\boldsymbol{U}+\boldsymbol{b}_{\sigma}-\boldsymbol{x_{0}}\|^{2}_{2}\\
 & =\mathbb{E}_{\boldsymbol{x_{0}}\sim p_{0},\boldsymbol{Z}\sim\mathcal{N}(0,\boldsymbol{I}),\boldsymbol{U}}\|\boldsymbol{W}^{\text{circ}}_{\sigma}(\boldsymbol{x_{0}}+\sigma\boldsymbol{Z})+\boldsymbol{V}_{\sigma}\boldsymbol{U}+(\boldsymbol{\boldsymbol{I}}-\boldsymbol{W}^{\text{circ}}_{\sigma})\boldsymbol{\mu_{p_{0}}}-\boldsymbol{V}_{\sigma}\boldsymbol{\mu_{U}}-\boldsymbol{x_{0}}\|^{2}_{2}\\
 & =\mathbb{E}_{\boldsymbol{x_{0}}\sim p_{0},\boldsymbol{Z}\sim\mathcal{N}(0,\boldsymbol{I}),\boldsymbol{U}}\|\boldsymbol{W}^{\text{circ}}_{\sigma}(\tilde{\boldsymbol{x_{0}}}+\sigma\boldsymbol{Z})-\tilde{\boldsymbol{x_{0}}}+\boldsymbol{V}_{\sigma}\tilde{\boldsymbol{U}}\|^{2}_{2}\\
 & =\text{Tr}\big(\boldsymbol{W}^{\text{circ}}_{\sigma}(\boldsymbol{\Sigma}_{\boldsymbol{p_{0}}}+\sigma^{2}\boldsymbol{I})\boldsymbol{W}^{\text{circ}\top}_{\sigma}-2\boldsymbol{W}^{\text{circ}}_{\sigma}\boldsymbol{\Sigma}_{\boldsymbol{p_{0}}}+\boldsymbol{\Sigma}_{\boldsymbol{p_{0}}}+\boldsymbol{V}_{\sigma}\boldsymbol{\Sigma}_{\boldsymbol{U}}\boldsymbol{V}^{\top}_{\sigma}+2\boldsymbol{V}_{\sigma}\text{Cov}(\boldsymbol{U},\boldsymbol{x_{0}})(\boldsymbol{W}^{\text{circ}}_{\sigma}-\boldsymbol{I})\big)\\
 & =\text{Tr}\big(\boldsymbol{W}^{\text{circ}}_{\sigma}(\boldsymbol{\Sigma}_{\boldsymbol{p_{0}}}+\sigma^{2}\boldsymbol{I})\boldsymbol{W}^{\text{circ}\top}_{\sigma}-2\boldsymbol{W}^{\text{circ}}_{\sigma}\boldsymbol{\Sigma}_{\boldsymbol{p_{0}}}+\boldsymbol{\Sigma}_{\boldsymbol{p_{0}}}-(\boldsymbol{\boldsymbol{I}-W}_{\sigma})\text{Cov}(\boldsymbol{x_{0}},\boldsymbol{U})\boldsymbol{\Sigma}^{-1}_{\boldsymbol{U}}\text{Cov}(\boldsymbol{U},\boldsymbol{x_{0}})(\boldsymbol{\boldsymbol{I}-W}_{\sigma})^{\top})\\
 & =\text{Tr}\big(\boldsymbol{W}^{\text{circ}}_{\sigma}(\boldsymbol{\Sigma}_{\boldsymbol{p_{0}}}+\sigma^{2}\boldsymbol{I}-\text{Cov}(\boldsymbol{x_{0}},\boldsymbol{U})\boldsymbol{\Sigma}^{-1}_{\boldsymbol{U}}\text{Cov}(\boldsymbol{U},\boldsymbol{x_{0}}))\boldsymbol{W}^{\text{circ}\top}_{\sigma}-2\boldsymbol{W}^{\text{circ}}_{\sigma}(\boldsymbol{\Sigma}_{\boldsymbol{p_{0}}}-\text{Cov}(\boldsymbol{x_{0}},\boldsymbol{U})\boldsymbol{\Sigma}^{-1}_{\boldsymbol{U}}\text{Cov}(\boldsymbol{U},\boldsymbol{x_{0}}))\\
 & +(\boldsymbol{\Sigma}_{\boldsymbol{p_{0}}}-\text{Cov}(\boldsymbol{x_{0}},\boldsymbol{U})\boldsymbol{\Sigma}^{-1}_{\boldsymbol{U}}\text{Cov}(\boldsymbol{U},\boldsymbol{x_{0}}))
\end{align*}
Let $\boldsymbol{\Sigma}_{\text{cond}}=\boldsymbol{\Sigma}_{\boldsymbol{p_{0}}}-\text{Cov}(\boldsymbol{x_{0}},\boldsymbol{U})\boldsymbol{\Sigma}^{-1}_{\boldsymbol{U}}\text{Cov}(\boldsymbol{U},\boldsymbol{x_{0}})$,
then:
\begin{align*}
\mathcal{L}^{\text{circ}}_{\sigma,\boldsymbol{U}} & =\text{Tr}\big(\boldsymbol{W}^{\text{circ}}_{\sigma}(\boldsymbol{\Sigma}_{\text{cond}}+\sigma^{2}\boldsymbol{I})\boldsymbol{W}^{\text{circ}\top}_{\sigma}-2\boldsymbol{W}^{\text{circ}}_{\sigma}\boldsymbol{\Sigma}_{\text{cond}}+\boldsymbol{\Sigma}_{\text{cond}}
\end{align*}
This is the same form as $\mathcal{L}^{\text{circ}}_{\sigma}$ except
with $\boldsymbol{\Sigma}_{\boldsymbol{p_{0}}}$ swapped out for $\boldsymbol{\Sigma}_{\text{cond}}$.
Hence we can draw the same conclusion: let $\boldsymbol{P}^{\text{cond}}=\boldsymbol{F}\boldsymbol{\Sigma}_{\text{cond}}\boldsymbol{F}^{*}$
and as before, $\text{diag}(\hat{\boldsymbol{W}})$ is the Discrete
Fourier Transform of the first row of the circulant matrix $\boldsymbol{W}$,
i.e. $\text{diag}(\hat{\boldsymbol{W}})_{j}\coloneqq\hat{w}_{j}=\sum^{d}_{k=0}w_{k}e^{-2\pi ijk/d}$
($w_{k}$ is the $k$-th entry of the first row of $\boldsymbol{W}$).
For $\boldsymbol{W}$ to have real entries, $\hat{w}_{j}=\bar{\hat{w}}_{d-j}$.
Let $\boldsymbol{P}^{\text{cond}}_{kk}=f^{H}_{k}\boldsymbol{\Sigma}_{\text{cond}}f_{k}$.
Then:
\[
\mathcal{L}^{\text{circ}}_{\sigma,\boldsymbol{U}}=\sum^{d}_{k=1}(\|\hat{w}_{k}\|^{2}(\boldsymbol{P}^{\text{cond}}_{kk}+\sigma^{2})-2\text{Re}(\hat{w}_{k})\boldsymbol{P}^{\text{cond}}_{kk})+\text{Tr}(\boldsymbol{\Sigma}_{\text{cond}})
\]
Recall that $\boldsymbol{\Sigma}_{\text{cond}}=\boldsymbol{\Sigma}_{\boldsymbol{p_{0}}}-\text{Cov}(\boldsymbol{x_{0}},\boldsymbol{U})\boldsymbol{\Sigma}^{-1}_{\boldsymbol{U}}\text{Cov}(\boldsymbol{U},\boldsymbol{x_{0}})=\text{Var}(\boldsymbol{x_{0}}\mid\boldsymbol{U})$
so $\boldsymbol{\Sigma}_{\text{cond}}$ is positive semi-definite
and so $\boldsymbol{P}^{\text{cond}}_{kk}$ is real and non-negative.
Similar to above, optimize for each $\hat{w}_{k}$ independently.
Let $\hat{w}_{k}=a+bi$ where $a,b\in\mathbb{R}$; we're minimizing:
\[
\|\hat{w}_{k}\|^{2}(\boldsymbol{P}^{\text{cond}}_{kk}+\sigma^{2})-2\text{Re}(\hat{w}_{k})\boldsymbol{P}^{\text{cond}}_{kk}=(a^{2}+b^{2})(\boldsymbol{P}^{\text{cond}}_{kk}+\sigma^{2})-2a\boldsymbol{P}^{\text{cond}}_{kk}
\]
Again it ought to be that $b=0$ and:
\[
a=\frac{\boldsymbol{P}^{\text{cond}}_{kk}}{\boldsymbol{P}^{\text{cond}}_{kk}+\sigma^{2}}=\hat{w}^{*}_{k}
\]
And so:
\[
\mathcal{L}^{\text{circ}}_{\sigma,\boldsymbol{U}}=\sigma^{2}\sum^{d}_{k=1}\frac{\boldsymbol{P}^{\text{cond}}_{kk}}{\boldsymbol{P}^{\text{cond}}_{kk}+\sigma^{2}}
\]
Recall that the optimal conditional linear denoiser loss without any
constraint on $\boldsymbol{W}_{\sigma}$ is:
\[
\mathcal{L}_{\sigma,\boldsymbol{U}}=\sigma^{2}\sum^{d}_{k=1}\frac{\lambda^{\text{cond}}_{k}}{\lambda^{\text{cond}}_{k}+\sigma^{2}}
\]
where $\lambda^{\text{cond}}_{k}$ is the eigenvalues of $\boldsymbol{\Sigma}_{\text{cond}}$,
and also the eigenvalues of $\boldsymbol{P}^{\text{cond}}$. Then
like before, by the Schur-Horn Theorem, for the Hermitian $\boldsymbol{P}^{\text{cond}}$,
$\boldsymbol{\lambda}^{\text{cond}}\succ\text{diag}(\boldsymbol{P}^{\text{cond}})$.
And so by the Karamata's Inequality/Majorization Theorem:
\[
\mathcal{L}^{\text{circ}}_{\sigma,\boldsymbol{U}}=\sigma^{2}\sum^{d}_{k=1}\frac{\boldsymbol{P}^{\text{cond}}_{kk}}{\boldsymbol{P}^{\text{cond}}_{kk}+\sigma^{2}}\geq\sigma^{2}\sum^{d}_{k=1}\frac{\lambda^{\text{cond}}_{k}}{\lambda^{\text{cond}}_{k}+\sigma^{2}}=\mathcal{L}_{\sigma,\boldsymbol{U}}
\]
Again in the conditional case, constraining $\boldsymbol{W}_{\sigma}$
to be circulant doesn't buy us anything more than the optimal linear
denoiser does. Also note that the gain from going from unconditional
to conditional in the case of circulant $\boldsymbol{W}_{\sigma}$
is the gain of going from $\text{Var}(\boldsymbol{x_{0}})=\boldsymbol{\Sigma}_{\boldsymbol{p_{0}}}$
to $\text{Var}(\boldsymbol{x_{0}}\mid\boldsymbol{U})=\boldsymbol{\Sigma}_{\text{cond}}$
w.r.t. getting the values of $\boldsymbol{P}_{kk}=f^{H}_{k}\boldsymbol{\Sigma}_{\boldsymbol{p_{0}}}f_{k}$
v.s. $\boldsymbol{P}^{\text{cond}}_{kk}=f^{H}_{k}\boldsymbol{\Sigma}_{\text{cond}}f_{k}$.

\section{Mutual Information of the Learned Distribution}

Recall that with parametrization $\gamma=1/t$, and $\gamma=1/\sigma^{2}$
then mutual information is:
\[
I(\boldsymbol{x_{0}};Y_{\gamma})=\frac{1}{2}\int^{\gamma}_{0}\text{MMSE}[\boldsymbol{x_{0}}\mid Y_{s}]\,ds
\]
Consider the affine circulant function class:
\[
\mathcal{F}_{\text{circ}}=\{y\mapsto\boldsymbol{W}_{\sigma}y+b:\boldsymbol{W}_{\sigma}\text{ is circulant}\}
\]
So the constrained MMSE by Tweedie's formula is:
\[
\text{MMSE}_{\mathcal{F}_{\text{circ}}}[\boldsymbol{x_{0}}\mid Y_{\gamma}]=\min_{\boldsymbol{W}\text{circulant},b}\mathbb{E}\|\boldsymbol{x_{0}}-(\boldsymbol{W}_{\sigma}Y_{\gamma}+b)\|^{2}_{2}
\]
From last section, with $\gamma=1/\sigma^{2}$, then:
\[
\text{MMSE}_{\mathcal{F}_{\text{circ}}}[\boldsymbol{x_{0}}\mid Y_{\gamma}]=\sum^{d}_{k=1}\frac{\boldsymbol{P}_{kk}}{1+\gamma\boldsymbol{P}_{kk}}
\]
Then the constrained mutual information is:
\begin{align*}
I_{\mathcal{F}_{\text{circ}}}(\boldsymbol{x_{0}};Y_{\gamma}) & =\frac{1}{2}\int^{\gamma}_{0}\sum^{d}_{k=1}\frac{\boldsymbol{P}_{kk}}{1+s\boldsymbol{P}_{kk}}\,ds\\
 & =\frac{1}{2}\sum^{d}_{k=1}\log(1+\gamma\boldsymbol{P}_{kk})
\end{align*}
Recall that in $\mathcal{F}_{L}$ the linear/affine space (unconstrained)
$\mathcal{F}_{L}=\{y\mapsto\boldsymbol{W}_{\sigma}y+b\}$, we have:
\[
I_{\mathcal{F}_{L}}(\boldsymbol{x_{0}};Y_{\gamma})=\frac{1}{2}\sum^{d}_{k=1}\log(1+\gamma\lambda_{k})
\]
where $\lambda_{k}$ is the eigenvalues of $\boldsymbol{\Sigma}_{\boldsymbol{p_{0}}}$.
Again we know that $\boldsymbol{\lambda}\succ\text{diag}(\boldsymbol{P})$
, and since $f(x)=\log(1+\gamma x)$ is concave, by the Karamata's
Inequality/Majorization Theorem, $\sum^{d}_{k=1}f(\lambda_{k})\leq\sum^{d}_{k=1}f(\boldsymbol{P}_{kk})$.
Therefore:
\[
I_{\boldsymbol{p_{0}}}(\boldsymbol{x_{0}};Y_{\gamma})\leq I_{\mathcal{F}_{L}}(\boldsymbol{x_{0}};Y_{\gamma})\leq I_{\mathcal{F}_{\text{circ}}}(\boldsymbol{x_{0}};Y_{\gamma})
\]
The second inequality also makes sense because 

On the other hand, entropy is:
\[
h(Y_{\gamma}\mid\boldsymbol{x_{0}})=\frac{d}{2}\log\Big(\frac{2\pi e}{\gamma}\Big)
\]
and
\[
h(Y_{\gamma})=\frac{d}{2}\log\Big(\frac{2\pi e}{\gamma}\Big)+I(\boldsymbol{x_{0}};Y_{\gamma})
\]
Therefore:
\begin{align*}
h(Y_{\gamma}) & \leq\frac{d}{2}\log\Big(\frac{2\pi e}{\gamma}\Big)+I_{\mathcal{F}_{\text{circ}}}(\boldsymbol{x_{0}};Y_{\gamma})\\
 & =\frac{1}{2}\sum^{d}_{k=1}\log\Big(2\pi e\big(\boldsymbol{P}_{kk}+\frac{1}{\gamma}\big)\Big)
\end{align*}
Again since $\gamma=1/\sigma^{2}$:
\[
h(\boldsymbol{x_{0}}+\sigma\boldsymbol{Z})\leq\frac{1}{2}\log\big(2\pi e(\boldsymbol{P}_{kk}+\sigma^{2})\big)
\]
For the conditional case:
\[
\text{MMSE}_{\mathcal{F}_{\text{circ}},U}[\boldsymbol{x_{0}}\mid Y_{\gamma},U]=\sum^{d}_{k=1}\frac{\boldsymbol{P}^{\text{cond}}_{kk}}{1+\gamma\boldsymbol{P}^{\text{cond}}_{kk}}
\]
So:
\begin{align*}
I_{\mathcal{F}_{\text{circ}}}(Y_{\gamma};U) & =\frac{1}{2}\int^{\gamma}_{0}\sum^{d}_{k=1}\text{MMSE}_{\mathcal{F}_{\text{circ}}}[\boldsymbol{x_{0}}\mid Y_{\gamma}]-\text{MMSE}_{\mathcal{F}_{\text{circ}},U}[\boldsymbol{x_{0}}\mid Y_{\gamma},U]\,ds\\
 & =\frac{1}{2}\sum^{d}_{k=1}\log\Big(\frac{1+\gamma\boldsymbol{P}_{kk}}{1+\gamma\boldsymbol{P}^{\text{cond}}_{kk}}\Big)
\end{align*}
Taking $\gamma\to\infty$:
\[
I_{\mathcal{F}_{\text{circ}}}(\boldsymbol{x_{0}};U)=\frac{1}{2}\sum^{d}_{k=1}\log\Big(\frac{\boldsymbol{P}_{kk}}{\boldsymbol{P}^{\text{cond}}_{kk}}\Big)
\]
Recall that in $\mathcal{F}_{L}$ the unconstrained linear/affine
space:
\[
I_{\mathcal{F}_{L}}(\boldsymbol{x_{0}};U)=\frac{1}{2}\sum^{d}_{k=1}\log\Big(\frac{\lambda_{k}}{\lambda^{\text{cond}}_{k}}\Big)
\]
Unlike the unconditional case we can't really compare $I_{\mathcal{F}_{\text{circ}}}(\boldsymbol{x_{0}};U)$
and $I_{\mathcal{F}_{L}}(\boldsymbol{x_{0}};U)$.

We can see that the circulant constraint increases its function class
MMSE and so increases the corresponding entropy upperbound.

\section{Unconditional Nonlinear Denoiser with Circulant Matrix}

Recall that the nonlinear map is: 
\[
\phi(\boldsymbol{y})=\omega(\Theta\boldsymbol{y}+\epsilon)
\]
where $\Theta\in\mathbb{R}^{k\times d}$ and $\epsilon\in\mathbb{R}^{k}$
are random matrices and the denoiser w.r.t this non-linear feature
map is:
\[
\boldsymbol{D}(\boldsymbol{y};\sigma)=\boldsymbol{W}_{\sigma}\phi(\boldsymbol{y})+\boldsymbol{b}_{\sigma}
\]
where $\boldsymbol{W}_{\sigma}\in\mathbb{R}^{d\times k}$ and $\boldsymbol{b}_{\sigma}\in\mathbb{R}^{d}$.
Since circulant matrices need to be square, assuming $k$ is divisible
by $d$ (which is the setup of our RF scaling experiment where we
drive $k/d$ to $\infty$), we can constrain $\Theta$ and $\boldsymbol{W}_{\sigma}$
to be circulant by block of $d\times d$. Let $k=c\cdot d$, write
$\Theta=\begin{bmatrix}\Theta_{1} & \Theta_{2} & \dots & \Theta_{c}\end{bmatrix}^{\top}$
where each $\Theta_{a}$ is a $d\times d$ circulant matrix. Correspondingly,
$\boldsymbol{W}_{\sigma}=\begin{bmatrix}\boldsymbol{W}_{1} & \boldsymbol{W}_{2} & \dots & \boldsymbol{W}_{c}\end{bmatrix}$,
each $\boldsymbol{W}_{a}$ for $a\in\{1,2,\dots,c\}$ is $d\times d$
circulant. Now, a circulant matrix is determined by the entries of
its first row, and the following rows after is just a cyclic shift
(each by one entry to the right), for example:
\[
\Theta_{a}=\begin{bmatrix}a_{1} & a_{2} & \dots & a_{d}\\
a_{d} & a_{1} & a_{2} & \dots\\
\vdots & \vdots & \vdots & \vdots\\
a_{2} & \dots & a_{d} & a_{1}
\end{bmatrix}
\]
Denote this cyclic shift as $\Theta_{a}=\text{circ}(h_{a})$ where
$h_{a}=\begin{bmatrix}a_{1} & a_{2} & \dots & a_{d}\end{bmatrix}$.
Under the random-feature framework, $h_{a}\sim\mathcal{N}(0,\frac{1}{d}\boldsymbol{I})$.
Hence each entry in the first row and so in the entire matrix has
marginal distribution $\mathcal{N}(0,1/d)$, the same as in the regular
case where there is no circulant constraint on $\Theta$ and $\boldsymbol{W}_{\sigma}$.
Our RF matrices in this case differ only in the joint distribution,
where each row in $\Theta_{a}$ is completely dependent on the first
row and each other; while in the regular case the rows are independent.
We have:
\[
\Theta\boldsymbol{y}=\begin{bmatrix}\Theta_{1}\boldsymbol{y} & \Theta_{2}\boldsymbol{y} & \dots & \Theta_{c}\boldsymbol{y}\end{bmatrix}^{\top}
\]
Let's index the rows if $\Theta$ with $\{j=(a,\tau)\}$ where $a\in{1,2,\dots,c}$
is indexing into the circulant blocks (there are $c$blocks) and $\tau=\{1,2,\dots,d-1\}$
is indexing into the rows of a circulant block starting from after
the first rows; each row is cyclic shift by $\tau$ entry to the right
of the first row $h_{a}=\begin{bmatrix}a_{1} & a_{2} & \dots & a_{d}\end{bmatrix}$,
i.e. $\Theta_{(a,\tau)}=S^{\tau}h_{a}$, for example $\Theta_{(a,1)}=S^{1}h_{a}=\begin{bmatrix}a_{d} & a_{1} & a_{2} & \dots\end{bmatrix}$,
$\Theta_{(a,d-1)}=S^{d-1}h_{a}=\begin{bmatrix}a_{2} & \dots & a_{d} & a_{1}\end{bmatrix}$ 

$S$ is the shift operator where $(Sx)_{i}=x_{i-1\mod d}$, so $(S^{\tau}x)_{i}=x_{i-\tau}$,
and $S^{\top}=S^{-1}$. 

For now let the noise $\epsilon$ be the same for each row in a block;
i.e. $\epsilon_{(a,\tau)}=\epsilon_{a}$, so noise is the same for
each feature in the same circulant block; since each row is a cyclic
shift of the first, this is to avoid random noise diluting the effect
of the cyclic shift (intuitively?, I'll come back to this later).
So $\epsilon$is drawn once and then fixed for each block. We have:
\[
[\Theta_{a}\boldsymbol{y}]_{\tau}=\Theta^{\top}_{(a,\tau)}\boldsymbol{y}=h^{\top}_{a}(S^{\tau})^{\top}\boldsymbol{y}=h^{\top}_{a}S^{-\tau}\boldsymbol{y}
\]
Let $\phi_{a}(\boldsymbol{y})=\omega(\Theta_{a}\boldsymbol{y}+\epsilon_{a})$,
then:
\[
\boldsymbol{D}(\boldsymbol{y};\sigma)=\sum^{c}_{a=1}\boldsymbol{W}_{a}\phi_{a}(\boldsymbol{y})+\boldsymbol{b}_{\sigma}
\]
Let each $\boldsymbol{W}_{a}=\boldsymbol{F}^{*}\text{diag}(\hat{\boldsymbol{W}_{a}})\boldsymbol{F}$,
where $\text{diag}(\hat{\boldsymbol{W}}_{a})_{j}\coloneqq\hat{w}^{a}_{j}=\sum^{d}_{k=0}w^{a}_{k}e^{-2\pi ijk/d}$
($w^{a}_{k}$ is the $k$-th entry of the first row of $\boldsymbol{W}_{a}$).
Recall from last write-up the denoiser loss is:
\[
\mathcal{L}^{\text{circ}}_{\sigma}=\text{Tr}(\boldsymbol{W}_{\sigma}\boldsymbol{\Sigma}_{\phi}\boldsymbol{W}^{\top}_{\sigma}-2\boldsymbol{W}_{\sigma}\boldsymbol{\Sigma}_{\phi,\boldsymbol{x_{0}}}+\boldsymbol{\Sigma}_{\boldsymbol{p_{0}}})
\]
$\boldsymbol{\Sigma}_{\phi}$ is $cd\times cd$, we can partition
into $c^{2}$ blocks of $d\times d$, denote each block $\boldsymbol{\Sigma}^{(a,b)}_{\phi}=\boldsymbol{\Sigma}_{\phi_{a},\phi_{b}}$
where $1\leq a,b\leq c$. $\boldsymbol{\Sigma}_{\phi,\boldsymbol{x_{0}}}$
is $cd\times d$, so we can partition into c blocks of $d\times d$
along its rows, denote each block $\boldsymbol{\Sigma}_{\phi_{a},\boldsymbol{x_{0}}}$.
Consider:
\[
\boldsymbol{W}_{\sigma}\boldsymbol{\Sigma}_{\phi}=\begin{bmatrix}\sum_{a}\boldsymbol{W}_{a}\boldsymbol{\Sigma}_{\phi_{a},\phi_{1}} & \sum_{a}\boldsymbol{W}_{a}\boldsymbol{\Sigma}_{\phi_{a},\phi_{2}} & \dots & \sum_{a}\boldsymbol{W}_{a}\boldsymbol{\Sigma}_{\phi_{a},\phi_{c}}\end{bmatrix}
\]
So:
\begin{align*}
\boldsymbol{W}_{\sigma}\boldsymbol{\Sigma}_{\phi}\boldsymbol{W}^{\top}_{\sigma} & =\sum_{a}\boldsymbol{W}_{a}\boldsymbol{\Sigma}_{\phi_{a},\phi_{1}}\boldsymbol{W}^{\top}_{1}+\sum_{a}\boldsymbol{W}_{a}\boldsymbol{\Sigma}_{\phi_{a},\phi_{2}}\boldsymbol{W}^{\top}_{2}+\dots+\sum_{a}\boldsymbol{W}_{a}\boldsymbol{\Sigma}_{\phi_{a},\phi_{d}}\boldsymbol{W}^{\top}_{d}\\
 & =\sum^{c}_{a=1}\sum^{c}_{b=1}\big(\boldsymbol{W}_{a}\boldsymbol{\Sigma}^{(a,b)}_{\phi}\boldsymbol{W}^{\top}_{b}\big)
\end{align*}
Use the fact that: 
\begin{align*}
\text{Tr}(\boldsymbol{W}_{a}\boldsymbol{M}\boldsymbol{W}^{\top}_{b}) & =\text{Tr}(\boldsymbol{F}^{*}\text{diag}(\hat{\boldsymbol{W}_{a}})\boldsymbol{F}\boldsymbol{M}\boldsymbol{F}^{*}\text{diag}(\hat{\boldsymbol{W}_{b}})^{H}\boldsymbol{F})\\
 & =\text{Tr}(\text{diag}(\hat{\boldsymbol{W}_{a}})\boldsymbol{F}\boldsymbol{M}\boldsymbol{F}^{*}\text{diag}(\hat{\boldsymbol{W}_{b}})^{H})
\end{align*}
then let $\boldsymbol{P}^{(a,b)}=\boldsymbol{F}\boldsymbol{\Sigma}^{(a,b)}_{\phi}\boldsymbol{F}^{*}$,
so $\boldsymbol{P}^{(a,b)}_{kk}=f^{H}_{k}\boldsymbol{\Sigma}^{(a,b)}_{\phi}f_{k}$
where $f_{k}$ is the $k$-th column of $\boldsymbol{F}^{*}$, or
the conjugate of the $k$-th row of $\boldsymbol{F}$. Note that $\boldsymbol{\Sigma}^{(a,b)}_{\phi}$
is not symmetric; rather $\boldsymbol{\Sigma}_{\phi}$ is symmetric
so $(\boldsymbol{\Sigma}^{(a,b)}_{\phi})^{\top}=\boldsymbol{\Sigma}^{(b,a)}_{\phi}$.
So here we can't apply the reasoning from the linear case for $a\neq b$
to make statements about $\boldsymbol{P}^{(a,b)}_{kk}$, so we separate
the pair indices into $(a,b)$ for $a\neq b$ and $(a,a)$. For $(a,a)$
indices, $\boldsymbol{P}^{(a,a)}_{kk}$ is real and non-negative,
and $\boldsymbol{P}^{(a,a)}_{kk}=\boldsymbol{P}^{(a,a)}_{d-k,d-k}$.
For $(a,b)$, we have:
\[
(\boldsymbol{P}^{(a,b)}_{kk})^{H}=f^{H}_{k}(\boldsymbol{\Sigma}^{(a,b)}_{\phi})^{H}(f^{H}_{k})^{H}=f^{H}_{k}\boldsymbol{\Sigma}^{(b,a)}_{\phi}f_{k}=\boldsymbol{P}^{(b,a)}_{kk}
\]
and since $f_{d-k}=\overline{f_{k}}$:
\begin{align*}
\boldsymbol{P}^{(a,b)}_{d-k,d-k} & =f^{H}_{d-k}\boldsymbol{\Sigma}^{(a,b)}_{\phi}f_{d-k}\\
 & =(\overline{f_{k}})^{H}\boldsymbol{\Sigma}^{(a,b)}_{\phi}\overline{f_{k}}=f^{\top}_{k}\boldsymbol{\Sigma}^{(a,b)}_{\phi}\overline{f_{k}}\\
 & =\text{conj}(f^{H}_{k}\boldsymbol{\Sigma}^{(a,b)}_{\phi}f_{k})=\text{conj}(\boldsymbol{P}^{(a,b)}_{kk})
\end{align*}
Consider the first term:
\begin{align*}
\text{Tr}(\boldsymbol{W}_{\sigma}\boldsymbol{\Sigma}_{\phi}\boldsymbol{W}^{\top}_{\sigma}) & =\sum^{c}_{a=1}\sum^{c}_{b=1}\text{Tr}(\text{diag}(\hat{\boldsymbol{W}_{a}})\boldsymbol{P}^{(a,b)}\text{diag}(\hat{\boldsymbol{W}_{b}})^{H})\\
 & =\sum^{c}_{a=1}\sum^{c}_{b=1}\sum^{d}_{k=1}\hat{w}^{a}_{k}\boldsymbol{P}^{(a,b)}_{kk}\text{conj}(\hat{w}^{b}_{k})
\end{align*}
Pair summand $(a,b)$ with $(b,a)$ for $a\neq b$, such that we can
cancel out the imaginary part of $\hat{w}^{a}_{k}$ and $\hat{w}^{b}_{k}$:
\begin{align*}
 & \hat{w}^{a}_{k}\boldsymbol{P}^{(a,b)}_{kk}\text{conj}(\hat{w}^{b}_{k})+\hat{w}^{b}_{k}\boldsymbol{P}^{(b,a)}_{kk}\text{conj}(\hat{w}^{a}_{k})\\
 & =\boldsymbol{P}^{(a,b)}_{kk}(\hat{w}^{a}_{k}\text{conj}(\hat{w}^{b}_{k}))+\boldsymbol{P}^{(b,a)}_{kk}\hat{w}^{b}_{k}\text{conj}(\hat{w}^{a}_{k})\\
 & =\boldsymbol{P}^{(a,b)}_{kk}(\hat{w}^{a}_{k}\text{conj}(\hat{w}^{b}_{k}))+\text{conj}(\boldsymbol{P}^{(a,b)}_{kk})\text{conj}(\hat{w}^{a}_{k}\text{conj}(\hat{w}^{b}_{k}))\\
 & =2\text{Re}(\hat{w}^{a}_{k}\boldsymbol{P}^{(a,b)}_{kk}\text{conj}(\hat{w}^{b}_{k}))
\end{align*}
Therefore:
\[
\text{Tr}(\boldsymbol{W}_{\sigma}\boldsymbol{\Sigma}_{\phi}\boldsymbol{W}^{\top}_{\sigma})=\sum^{c}_{a\neq b}\sum^{d}_{k=1}2\text{Re}(\hat{w}^{a}_{k}\boldsymbol{P}^{(a,b)}_{kk}\text{conj}(\hat{w}^{b}_{k}))+\sum^{c}_{a=1}\sum^{d}_{k=1}\|\hat{w}^{a}_{k}\|^{2}\boldsymbol{P}^{(a,a)}_{kk}
\]
Next, we evaluate the cross-covariance term $-2\text{Tr}(W_{\sigma}\Sigma_{\phi,x_{0}})$.
Partitioning $\Sigma_{\phi,x_{0}}$ into $c$ blocks of $d\times d$
along its rows, denoted as $\Sigma_{\phi_{a},x_{0}}$, then:
\[
\boldsymbol{W}_{\sigma}\Sigma_{\phi,\boldsymbol{x_{0}}}=\sum^{c}_{a=1}\boldsymbol{W}_{a}\Sigma_{\phi_{a},\boldsymbol{x_{0}}}
\]
and
\[
\text{Tr}(\boldsymbol{W}_{\sigma}\Sigma_{\phi,\boldsymbol{x_{0}}})=\sum^{c}_{a=1}\text{Tr}(\text{diag}(\hat{\boldsymbol{W}}_{a})\boldsymbol{F}\Sigma_{\phi_{a},\boldsymbol{x_{0}}}\boldsymbol{F}^{*})
\]
Let $\boldsymbol{Q}^{(a)}=\boldsymbol{F}\Sigma_{\phi_{a},\boldsymbol{x_{0}}}\boldsymbol{F}^{*}$
with diagonal entries $\boldsymbol{Q}^{(a)}_{kk}=f^{H}_{k}\boldsymbol{\Sigma}_{\phi_{a},\boldsymbol{x_{0}}}f_{k}$.
Then:
\[
\text{Tr}(\boldsymbol{W}_{\sigma}\boldsymbol{\Sigma}_{\phi,\boldsymbol{x_{0}}})=\sum^{c}_{a=1}\sum^{d}_{k=1}\hat{w}^{a}_{k}\boldsymbol{Q}^{(a)}_{kk}
\]
Since $\boldsymbol{\Sigma}_{\phi_{a},\boldsymbol{x_{0}}}$ has real
entries, $\boldsymbol{Q}^{(a)}_{d-k,d-k}=\text{conj}(\boldsymbol{Q}^{(a)}_{kk})$.
Also since $\boldsymbol{W}_{\sigma}$ has real entries, $\hat{w}^{a}_{j}=\text{conj}(\hat{w}^{a}_{d-j})$,
pairing $k$ and $d-k$ gives: 
\[
\hat{w}^{a}_{k}\boldsymbol{Q}^{(a)}_{kk}+\hat{w}^{a}_{d-k}\boldsymbol{Q}^{(a)}_{d-k,d-k}=2\text{Re}(\hat{w}^{a}_{k}\boldsymbol{Q}^{(a)}_{kk})
\]
Plugging back into the loss expressions:
\[
\mathcal{L}^{\text{circ}}_{\sigma}=\sum^{d}_{k=1}\left[\sum^{c}_{a=1}\|\hat{w}^{a}_{k}\|^{2}\boldsymbol{P}^{(a,a)}_{kk}+\sum^{c}_{a\ne b}2\text{Re}(\hat{w}^{a}_{k}\boldsymbol{P}^{(a,b)}_{kk}\text{conj}(\hat{w}^{b}_{k}))-2\sum^{c}_{a=1}\text{Re}(\hat{w}^{a}_{k}\boldsymbol{Q}^{(a)}_{kk})\right]+\text{Tr}(\boldsymbol{\Sigma}_{\boldsymbol{p_{0}}})
\]
To optimize, we group the complex weights of all $c$ blocks into
a vector $\hat{\mathbf{w}}_{k}=[\hat{w}^{1}_{k},\dots,\hat{w}^{c}_{k}]^{\top}\in\mathbb{C}^{c}$.
Define the Hermitian matrix $\mathbf{P}_{k}\in\mathbb{C}^{c\times c}$
where $(\mathbf{P}_{k})_{ab}=\boldsymbol{P}^{(a,b)}_{kk}$, and the
vector $\mathbf{q}_{k}\in\mathbb{C}^{c}$ where $(\mathbf{q}_{k})_{a}=\boldsymbol{Q}^{(a)}_{kk}$.
Let $\boldsymbol{v}_{k}\coloneqq\text{conj}(\hat{\mathbf{w}}_{k})$.
Then $\sum^{c}_{a\ne b}\hat{w}^{a}_{k}\boldsymbol{P}^{(a,b)}_{kk}\text{conj}(\hat{w}^{b}_{k})=\boldsymbol{v}^{H}_{k}\mathbf{P}_{k}\boldsymbol{v}_{k}$
and $\sum^{c}_{a=1}\hat{w}^{a}_{k}\boldsymbol{Q}^{(a)}_{kk}=\boldsymbol{v}^{H}_{k}\mathbf{q}_{k}$.
Similar to above it looks like we can optimize w.r.t to each index
$k$ independently, however note that there is one constraint $\hat{\mathbf{w}}_{k}=\text{conj}(\hat{\mathbf{w}}_{d-k})$,
but since $\mathbf{P}_{k}=\text{conj}(\mathbf{P}_{d-k})$, optimizing
for both independently is still valid since the two independent solution
will automatically satisfy this constraint. The objective for a given
$k$ is the quadratic form:
\begin{align*}
\mathcal{L}^{\text{circ}}_{\sigma,k} & =\boldsymbol{v}^{H}_{k}\mathbf{P}_{k}\boldsymbol{v}_{k}-2\text{Re}(\boldsymbol{v}^{H}_{k}\mathbf{q}_{k})\\
 & =\boldsymbol{v}^{H}_{k}\mathbf{P}_{k}\boldsymbol{v}_{k}-\boldsymbol{v}^{H}_{k}\mathbf{q}_{k}-\mathbf{q}^{H}_{k}\boldsymbol{v}_{k}
\end{align*}
Taking the Wirtinger derivative w.r.t. $\boldsymbol{v}^{H}_{k}$ and
setting it to $0$:

\[
\mathbf{P}_{k}\boldsymbol{v}_{k}-\mathbf{q}_{k}=0
\]
so
\[
\boldsymbol{v}^{*}_{k}=\mathbf{P}^{-1}_{k}\mathbf{q}_{k}
\]
Plugging back to get the loss:
\[
\mathcal{L}^{\text{circ}}_{\sigma}=\text{Tr}(\boldsymbol{\Sigma}_{\boldsymbol{p_{0}}})-\sum^{d}_{k=1}\mathbf{q}^{H}_{k}\mathbf{P}^{-1}_{k}\mathbf{q}_{k}
\]
Now it might be misleading to try to compare $\mathcal{L}^{\text{circ}}_{\sigma}$
with $\mathcal{L}_{\sigma}=\text{Tr}(\boldsymbol{\Sigma}_{\boldsymbol{p_{0}}}-\text{Cov}(\boldsymbol{x_{0}},\phi)\boldsymbol{\Sigma}^{-1}_{\phi}\text{Cov}(\phi,\boldsymbol{x_{0}}))$
since $\mathcal{L}_{\sigma}$ is the unconstrained denoiser with dense
RF $\Theta$ inside $\phi$, however recall that $\Theta^{\text{circ}}$
is not a subset of $\Theta^{\text{dense}}$ since the joint distributions
of the rows are different: rows of $\Theta^{\text{circ}}$ are dependent
while those of $\Theta^{\text{dense}}$ are independent. Thus we can
no longer make a blanket claim that $\mathcal{L}^{\text{circ}}_{\sigma}\geq\mathcal{L}_{\sigma}$
always, even though the $\boldsymbol{W}^{\text{circ}}_{\sigma}\subset\boldsymbol{W}_{\sigma}$
still. The point is that given the same $\Theta$, since $\boldsymbol{W}^{\text{circ}}_{\sigma}$
and $\boldsymbol{W}_{\sigma}$ are trainable, the loss for the denoiser
with the contrained $\boldsymbol{W}^{\text{circ}}_{\sigma}$ can never
be more optimal than that of the unconstrained, optimal $\boldsymbol{W}_{\sigma}$.
However $\Theta$ in our circulant setting and $\Theta$ in the regular
unconstrained nonlinear denoiser is not the same; thus leaving space
for $\mathcal{L}^{\text{circ}}_{\sigma}$ to beat $\mathcal{L}_{\sigma}$

\section{When does $\mathcal{L}^{\text{circ}}_{\sigma}$ beat $\mathcal{L}_{\sigma}$
in nonlinear denoiser?}

Through our experiments the RF denoiser needs a lot of features, i.e.
$k/d\to\infty$. We know that circulant structure gives its best loss
when the data is stationary, for example in a face dataset, the nose
can equally be in the pixel at the top of the image as the pixel in
the middle. We want to know if the circulant structure can help when
we have finite $k$. To formalize this we first decompose the loss
between the circulant structure and the dense structure:
\[
\mathcal{L}^{\text{circ}}(k)-\mathcal{L}^{\text{dense}}(k)=A_{\text{circ}}(k)+\Delta_{\text{stat}}(\sigma)-A_{\text{dense}}(k)
\]
where 
\begin{align*}
A_{\text{dense}}(k) & :=\mathcal{L}^{\text{dense}}(k)-\text{MMSE}(\boldsymbol{p_{0}})\\
A_{\text{circ}}(k) & \coloneqq\mathcal{L}^{\text{circ}}(k)-\text{MMSE}(\bar{\boldsymbol{p_{0}}})\\
\bar{\boldsymbol{p_{0}}}(x) & =\frac{1}{d}\sum^{d-1}_{\tau=0}\boldsymbol{p_{0}}(S^{-\tau}x)\\
\ensuremath{\Delta_{\text{stat}}(\sigma)} & :=\text{MMSE}(\bar{\boldsymbol{p_{0}}})-\text{MMSE}(\boldsymbol{p_{0}})
\end{align*}
is a stationary/shift-variant distribution made from shift-symmetrizing
the true data distribution i.e. $\boldsymbol{p_{0}}$, by applying
a cyclic shift operator $S^{-\tau}$ to every possible shift $\tau$,
and then averaging the result.

By constraining the weights to a circulant matrix, the network is
forced to be shift-equivariant. Therefore, when evaluating the non-stationary
data distribution $\boldsymbol{p_{0}}$, the constrained denoiser
mathematically optimizes as if the data were drawn from the shift-symmetrized
(stationary) distribution $\bar{\boldsymbol{p_{0}}}$. Theoretically
it can never be better than $\text{MMSE}(\bar{\boldsymbol{p_{0}}})$.
This cost is quantified by $\ensuremath{\Delta_{\text{stat}}(\sigma)}$.
Note that $\bar{\boldsymbol{p_{0}}}$ is generatively $\bar{\boldsymbol{x_{0}}}=S^{T}\boldsymbol{x_{0}}$,
where the shift is $T\sim\text{Unif}\{0,\dots,d-1\}$. Because conditioning
on $T$ brings the problem back to the unshifted, non-stationary $\boldsymbol{p_{0}}$,
the stationarization penalty is equivalent to the gap between unconditional
and conditional denoising:

\[
\Delta_{\text{stat}}(\sigma)=\text{MMSE}_{\bar{\boldsymbol{p_{0}}}}[\boldsymbol{X}\mid\boldsymbol{Y}_{\sigma}]-\text{MMSE}_{\bar{\boldsymbol{p_{0}}}}[\boldsymbol{X}\mid\boldsymbol{Y}_{\sigma},T]
\]
By the I-MMSE identity, integrating this penalty over the noise schedule
gives the mutual information between the shifted data and the shift
operator, bounded by the entropy of $T$:

\[
\frac{1}{2}\int^{\infty}_{0}\frac{dt}{t^{2}}\Delta_{\text{stat}}(t)=I(\bar{\boldsymbol{x_{0}}};T)\le H(T)=\log d
\]
Because the penalty $\Delta_{\text{stat}}$ is bounded by $\log d$
across the noise schedule, it is a finite constant independent of
$k$. Conversely, at small $k$, $A_{\text{dense}}(k)$ is unconstrained/unbounded
due to the $O(d^{n})$ parameters required to cover the dense feature
space. Furthermore, $A_{\text{circ}}$ is a stationary problem on
$\bar{\boldsymbol{p_{0}}}$, therefore we can derive this as it decomposes
per-frequency like in the last section. $\boldsymbol{W}=0$ is possible,
so $\mathcal{L}^{\text{circ}}\leq\text{Tr}(\boldsymbol{\Sigma}_{\boldsymbol{p_{0}}})$
(mean-predicting), so $0\leq A_{\text{circ}}(k)\leq\text{Tr}(\boldsymbol{\Sigma}_{\boldsymbol{p_{0}}})-\text{MMSE}(\bar{\boldsymbol{p_{0}}})$.
If we consider the difference:
\[
\mathcal{L}^{\text{circ}}(k)-\mathcal{L}^{\text{dense}}(k)=A_{\text{circ}}(k)+\Delta_{\text{stat}}(\sigma)-A_{\text{dense}}(k)
\]

If $k\to0$, both models degenerate to the bias-only predictor, $\mathcal{L}^{\text{circ}},\mathcal{L}^{\text{dense}}\to\text{Tr}(\boldsymbol{\Sigma}_{\boldsymbol{p_{0}}})$,
so $A_{\text{dense}}(k)-A_{\text{circ}}(k)\to\Delta_{\text{stat}}$,
so the difference $\to0$. On the other hand if $k\to\infty$, both
$A$ vanishes, and the difference $\to\Delta_{\text{stat}}>0$, so
dense would win.

So the win region is where $A_{\text{dense}}(k)-A_{\text{circ}}(k)$
rises above its starting value $\Delta_{\text{stat}}$ before decaying
to $0$ along $k$, which can be non-monotone. Point is if we can
compute these quantities from our data set then we get a set of $k$
where the circulant denoiser outperforms the dense denoiser.
\end{document}
